\documentclass[10pt,a4paper]{amsart}
\usepackage[margin=19mm]{geometry}
\usepackage{amsmath,amssymb,mathtools}
\usepackage[scr=rsfs,cal=euler]{mathalfa}
\usepackage{xfrac}
\usepackage[unicode,hidelinks,bookmarks=false]{hyperref}
\newcommand{\ie}{i.e.\ }
\newcommand{\cf}{cf.\ }
\newcommand{\resp}{respectively\ }
\newcommand{\Subsection}{\subsection}
\providecommand{\nobreakdash}{\nobreak\hbox{-}}
\setlength{\emergencystretch}{2em}
\multlinegap0pt
\usepackage{tikz}
\usetikzlibrary{cd}
\usepackage{amssymb}
\usepackage{mathtools}
\usepackage{tikz}
\usepackage{tcolorbox}
\newtheorem{prop}{Proposition}
\newcommand{\End}{\operatorname{End}}
\newcommand{\Hom}{\operatorname{Hom}}
\newcommand{\Ker}{\operatorname{Ker}}
\def\La{\Lambda}
\def\bbN{\mathbb N}
\def\bfL{\mathbf L}
\newcommand{\colim}{\operatorname*{colim}}
\renewcommand{\mod}{\operatorname{mod}}
\title{From finite to infinite length modules\\ over tame hereditary algebras}
\author{Lidia Angeleri H\"ugel, Andrew Hubery, Henning Krause}
\date{Source: arXiv:2604.27528v2 (CC BY 4.0)}
\begin{document}
\maketitle

\noindent\emph{Source and licence.} From finite to infinite length modules\\ over tame hereditary algebras. Version arXiv:2604.27528v2 (CC BY 4.0), distributed by arXiv under the Creative Commons Attribution 4.0 International licence, http://creativecommons.org/licenses/by/4.0/, as recorded in the arXiv licence metadata for this exact version and shown on the abstract page https://arxiv.org/abs/2604.27528v2. Full licence notice: \texttt{licenses/arxiv-2604.27528-CC-BY-4.0.txt}.\par\smallskip

\noindent\textit{Editorial scope.} Counted unit: the article's prop pr:subgroup-fin-def (source0.tex lines 675--680). The article's complete original proof of it is delivered with it (source0.tex lines 682--721, one \texttt{proof} environment of 40 lines). No mathematics is written, repaired or re-derived: the statement and the proof are the article's own lines, in the article's own notation, and no retained line is substituted or re-flowed.  No further source-local context is needed: every cross-reference of every retained line resolves inside the two delivered spans.  Nothing was omitted from any retained span, so the package carries no omission record.  No retained line contains a citation, so the wrapper carries no bibliography and no bibliography entry.  The retained lines contain the article's own diagram code, which the wrapper typesets with the standard tikz packages; no image file is delivered and the package declares no asset dependency.  Editorial formatting only, in addition: an AMS \texttt{amsart} preamble that restates the article's own declarations and macros used by the retained lines, namely the environments \texttt{prop} on separate counters, together with the standard packages those lines need.  The retained environments print in this document as Proposition 1 in order of appearance, whereas the article prints the same items with the numbering of its own full document.  The complete native source of the article as distributed in the arXiv e-print for this exact version is bundled unchanged as \texttt{sources/199479/source0.tex}.
\begin{prop}\label{pr:subgroup-fin-def}
\begin{enumerate}
\item A module is $\Sigma$-pure injective if and only if it has the descending chain condition (dcc) on subgroups of finite definition.
\item A module $X$ is endofinite if and only if it has both acc and dcc on subgroups of finite definition, in which case $\bfL_C(X)$ identifies with the lattice of submodules of $\Hom_\Lambda(C,X)$.
\end{enumerate}
\end{prop}

\begin{proof}
  We begin by observing that the chain conditions hold for $\bfL(X)$
  if and only if they hold for all $\bfL_C(X)$. One direction is
  clear, taking $C=\Lambda$. For the converse observe that any
  coherent functor $G$ is a subquotient of a finite direct sum of
  copies of the forgetful functor $F=\Hom_\La(\La,-)$.
  Thus it follows from the above description of lattices of subgroups
  of finite definition given by extensions of coherent functors that any chain
  condition for $\bfL(X)=\bfL_F(X)$ implies the condition for
  $\bfL_G(X)$. It remains to note that  $\bfL_C(X)=\bfL_G(X)$ for $G=\Hom_\La(C,-)$. 

(1) Assume $X$ has dcc on subgroups of definition. Then so too does $X^{(I)}$, since $\bfL(X^{(I)})=\bfL(X)$, and it is enough to show that $X$ is pure-injective.

Take a pure exact sequence $0\to X\to Y\to Z\to 0$ and write $Z=\colim Z_\alpha$ with $Z_\alpha\in\mod\Lambda$. This yields the inverse system of short exact sequences
\[ 0 \to \Hom_\Lambda(Z_\alpha,X) \to \Hom_\Lambda(Z_\alpha,Y) \to \Hom_\Lambda(Z_\alpha,Z) \to 0. \]
Moreover, the condition on $\bfL_{Z_\alpha}(X)$ shows that the Mittag-Leffler condition holds, so taking limits yields the exact sequence
\[ 0 \to \Hom_\Lambda(Z,X) \to \Hom_\Lambda(Z,Y) \to \Hom_\Lambda(Z,Z) \to 0. \]
It follows that the original sequence splits, and hence that $X$ is pure-injective.

Conversely assume that $X$ is $\Sigma$-pure injective. Then the canonical inclusion $X^{(\bbN)}\to X^\bbN$ admits a left inverse $\pi$. Assume for contradiction that we have a strict chain $\Hom_\Lambda(C,X)=U_0\supset U_1\supset\cdots$ in $\bfL_C(X)$ and for each $i$ take $f_i\in U_i-U_{i+1}$. This yields an element $f=(f_i)\in\Hom_\Lambda(C,X^\bbN)$. Take $n$ such that $(\pi f)_r=0$ for all $r\geq n$, and decompose $f=f'+f''$ having supports $\leq n$ and $>n$ respectively. Then $\pi f=f'+\pi f''$, so $f_n=f'_n=-(\pi f'')_n$.

Now, $f''$ lies in $U_{n+1}^\bbN$, which is a subgroup of finite definition in $\Hom_\Lambda(C,X^\bbN)$, and hence an $\End_\Lambda(X^\bbN)$-submodule. Thus $\pi f''$ must lie in $U_{n+1}^\bbN$, so $f_n=-(\pi f'')_n$ lies in $U_{n+1}$, a contradiction.
\medskip

(2) As all subgroups of finite definition of $X$ are $\End_\Lambda(X)$-submodules, we see that $X$ being endofinite implies $\bfL_C(X)$ has both acc and dcc.

For the converse assume first that $X$ is pure-injective. Write $X=\colim X_\alpha$ with $X_\alpha\in\mod\Lambda$ and associated morphisms $j_\alpha\colon X_\alpha\to X$. Given $\phi\colon C\to X$ with $C\in\mod\Lambda$, we can find $\beta$ such that $\phi=j_\beta\phi_\beta$. We claim that, inside $\Hom_\Lambda(C,X)$, we have
\[ \End_\Lambda(X)\phi = \bigcap_{\alpha\geq\beta}\Hom_\Lambda(X_\beta,X)\phi_\alpha. \]
One inclusion is clear, so suppose $\psi$ lies in the intersection. We factor $\phi=\bar\phi\pi$ and $\phi_\alpha=\bar\phi_\alpha\pi_\alpha$ via their respective images $\bar C$ and $\bar C_\alpha$, yielding commutative diagrams
\[ \begin{tikzcd}
\varepsilon_\alpha \colon &[-20pt] 0 \arrow[r] & \bar C_\alpha \arrow[r,"\bar\phi_\alpha"] \arrow[d,"i_\alpha"] &
X_\alpha \arrow[r] \arrow[d,"j_\alpha"] & Y_\alpha \arrow[r] \arrow[d] & 0\\
\varepsilon \colon &[-20pt] 0 \arrow[r] & \bar C \arrow[r,"\bar\phi"] & X \arrow[r] & Y \arrow[r] & 0
\end{tikzcd} \]
Moreover, $\varepsilon=\colim\varepsilon_\alpha$, which agrees with the colimit of the pushout sequences $\colim i_\alpha\varepsilon_\alpha$.

By assumption we have $\psi=\psi_\alpha\bar\phi_\alpha\pi_\alpha$. Using $\pi=i_\alpha\pi_\alpha$ we see that $\Ker\pi=\colim\Ker\pi_\alpha$, and hence $\psi$ factors as $\tilde\psi\pi$. We deduce that $\psi_\alpha\bar\phi_\alpha=\bar\psi i_\alpha$. Finally, the pushout $\bar\psi\varepsilon$ is the colimit of $\bar\psi i_\alpha\varepsilon_\alpha=\psi_\alpha\bar\phi_\alpha\varepsilon_\alpha$. The latter are all split exact, so $\bar\psi\varepsilon$ is pure, and hence also split as $X$ is pure-injective. Thus $\bar\psi=\gamma\bar\phi$ for some $\gamma\in\End_\Lambda(X)$, and $\psi=\gamma\phi$ factors through $\phi$.

If now $X$ is $\Sigma$-pure injective, then it has dcc on subgroups of finite definition, so each cyclic $\End_\Lambda(X)$-submodule of $\Hom_\Lambda(C,X)$ is a subgroup of finite definition. If we also assume acc, then the same holds for every $\End_\Lambda(X)$-submodule, showing that $X$ is endofinite.
\end{proof}

\end{document}
